\chapter{时空过程、相对论火箭与延迟成像}
\section{坐标原点与仿射变换}\label{th:m2affine}
洛伦兹变换作用于事件差。若两系并不把同一个事件取为零点，坐标关系应写为$x'^\mu=\Lambda^\mu{}_\nu x^\nu+C^\mu$。常数$C^\mu$由一件已知事件的坐标确定。平移不改变速度、微分间隔和固有时间，却会改变某个事件的坐标读数。

例如一只钟在事件$B=(ct_B,x_B)$换为速度$-V$匀速运动，并规定它在$B$显示$\tau_B$。以该钟为$x_2=0$原点的返程系可以定义为
\begin{align}\label{eq:m2return}
 ct_2&=c\tau_B+\gamma[(ct-ct_B)+\beta(x-x_B)],\\
 x_2&=\gamma[(x-x_B)+\beta(ct-ct_B)].\notag
\end{align}
把$B$代入，得到$(ct_2,x_2)=(c\tau_B,0)$；令$x_2=0$，又得到$x-x_B=-V(t-t_B)$，所以原点世界线确实沿返程运动。这两个代入可以避免把平移常数的符号写反。

\begin{example}\label{ex:m2twins}
地球与恒星在$S$中相距$l=10$光年。旅行者由地球以$V=0.8c$飞往恒星，立即换向并以同速返回。忽略转向所占时间。求地球和旅行者经过的时间；建立去程、返程坐标，并求转向前后旅行者判定与转向同时的地球时刻。
\Diagram{fig:m2twinsq}{地球与旅行者的世界线及转向事件}{\draw[axis](-.4,0)--(4,0)node[right]{$x$};\draw[axis](0,-.2)--(0,5.8)node[above]{$ct$};\draw[curve](0,0)--(0,5);\draw[curve](0,0)--(2,2.5)--(0,5);\draw[guide](2,0)--(2,5.2);\node[below]at(2,0){$l$};\fill(2,2.5)circle(1.4pt)node[right]{$B$};\node[left]at(0,5){$C$};\node[below left]at(0,0){$A$};}
\end{example}
\begin{solution}
图\ref{fig:m2twinsq}中$B$的坐标为$(cl/V,l)$，$C$的坐标为$(2cl/V,0)$。地球时间$T=2l/V=25$年。每段旅行者时间为$l/(\gamma V)$，$\gamma=5/3$，故总固有时间15年，去返各7.5年。

去程系原点由$A$出发，使用$ct_1=\gamma(ct-\beta x)$、$x_1=\gamma(x-\beta ct)$。在$B$，$ct_{1B}=l/(\gamma\beta)$。返程原点必须经过$B$且沿$-V$运动，钟面连续，故式\eqref{eq:m2return}取$ct_B=l/\beta$、$x_B=l$、$c\tau_B=l/(\gamma\beta)$，整理得
\[
 ct_2=\gamma(ct+\beta x)-2\gamma\beta l,
 \qquad x_2=\gamma(x+\beta ct)-2\gamma l.
\]
逆变换为
\[
 ct=\gamma(ct_2-\beta x_2),\qquad
 x=\gamma(x_2-\beta ct_2)+2l.
\]
\Diagram{fig:m2twinsa}{转向事件处的两条同时线与地球世界线相交于$D_-,D_+$}{\draw[axis](-.2,0)--(4,0)node[right]{$x$};\draw[axis](0,-.2)--(0,5.7)node[above]{$ct$};\draw[curve](0,0)--(2,2.5)--(0,5);\draw[guide](0,.9)--(3.1,3.38);\draw[guide](0,4.1)--(3.1,1.62);\fill(2,2.5)circle(1.4pt)node[right]{$B$};\node[left]at(0,.9){$D_-$};\node[left]at(0,4.1){$D_+$};}
去程同时线满足$ct-\beta x=ct_B-\beta l$，与地球$x=0$相交得
\[
 t_- =\frac{l}{\beta c}-\frac{\beta l}c=4.5\text{年}.
\]
返程同时线满足$ct+\beta x=ct_B+\beta l$，给$t_+=l/(\beta c)+\beta l/c=20.5$年。图\ref{fig:m2twinsa}中两交点相差$2\beta l/c=16$年。这是改变同时性定义后对远处地球时刻的重新赋值，地球钟本身连续走时，旅行者在转向处也没有瞬间看到地球老16年。
\end{solution}

\begin{example}\label{ex:m2fixed}
在例\ref{ex:m2twins}的坐标约定下，求同时满足$ct_2=ct$、$x_2=x$的事件。说明它是否位于旅行者转向处。
\Diagram{fig:m2fixedq}{返程系的坐标网原点世界线经过$B$}{\draw[axis](-.3,0)--(4,0)node[right]{$x$};\draw[axis](0,-.2)--(0,5.6)node[above]{$ct$};\draw[curve](0,0)--(2,2.5)--(0,5);\node[right]at(2,2.5){$B$};\draw[guide](2,0)--(2,5.2);}
\end{example}
\begin{solution}
将相同坐标条件代入逆变换，得到
\[
 (\gamma-1)ct=\gamma\beta x,
 \qquad(\gamma-1)x-\gamma\beta ct=-2l.
\]
第一式给$ct=\gamma\beta x/(\gamma-1)$。代入第二式后，$x$的系数为
\[
 (\gamma-1)-\frac{\gamma^2\beta^2}{\gamma-1}
 =(\gamma-1)-(\gamma+1)=-2,
\]
所以$x=l$，$ct=l(\gamma+1)/(\gamma\beta)$。取本题参数得到$ct=2l$，即$t=20$年。
\Diagram{fig:m2fixeda}{坐标不动事件$F$与转向事件$B$不同}{\draw[axis](-.2,0)--(3.8,0)node[right]{$x/l$};\draw[axis](0,-.2)--(0,4.6)node[above]{$ct/l$};\draw[curve](0,0)--(1.5,1.875)--(0,3.75);\fill(1.5,1.875)circle(1.4pt)node[right]{$B:(1.25,1)$};\fill(1.5,3)circle(1.4pt)node[right]{$F:(2,1)$};\draw[guide](1.5,0)--(1.5,3.8);}
图\ref{fig:m2fixeda}中$B$的地球时刻12.5年，$F$是恒星处更晚的事件。不同坐标系给某件事相同的数字，与该事件处是否有两物体相遇无关；此结果依赖既定的坐标平移约定。
\end{solution}

\section{共线变换与快度}\label{th:m2rapidity}
沿同一方向的洛伦兹矩阵在$(ct,x)$子空间内为
\[
 \Lambda(\beta)=\begin{pmatrix}\gamma&-\gamma\beta\\-\gamma\beta&\gamma\end{pmatrix}.
\]
将两矩阵相乘，其非对角元与对角元之比给
\[
 \beta_{12}=\frac{\beta_1+\beta_2}{1+\beta_1\beta_2},\quad
 \gamma_{12}=\gamma_1\gamma_2(1+\beta_1\beta_2).
\]
定义快度$\chi$满足$\beta=\tanh\chi$，于是$\gamma=\cosh\chi$、$\gamma\beta=\sinh\chi$。双曲函数加法公式说明共线快度直接相加。若使用虚数时间$ict$，形式上可把双曲变换写成虚角转动；这里的角度只是代数类比，不表示在普通空间中转动物体。

对一维运动，$u^\mu=c(\cosh\chi,\sinh\chi)$，故
\[
 a^\mu=c\frac{\dd\chi}{\dd\tau}(\sinh\chi,\cosh\chi),\quad
 \alpha=\sqrt{-a^2}=c\left|\frac{\dd\chi}{\dd\tau}\right|.
\]
恒定固有加速度对应快度随固有时线性增加。实验室速度则按双曲正切变化。

\begin{example}\label{ex:m2boosts}
飞船相对地球以$0.6c$前进，在船上又沿前方发射相对船速$0.6c$的探测器。求探测器相对地球的速度、快度和洛伦兹因子。
\Diagram{fig:m2boostsq}{两次同向速度变换}{\draw[axis](-.3,0)--(6,0)node[right]{$x$};\draw[curve](0,.3)rectangle(2,1);\node[above]at(1,1){飞船};\draw[axis](2.2,.65)--(3.4,.65)node[above]{$0.6c$};\fill(4.3,.65)circle(2pt);\draw[axis](4.4,.65)--(5.6,.65)node[above]{$0.6c$船系};}
\end{example}
\begin{solution}
图\ref{fig:m2boostsq}的两次速度属于不同参考系。令飞船和探测器在各自上一级参考系中的速度为$u_1=u_2=.6c$，实验室速度为
\[
 \beta=\frac{.6+.6}{1+.6^2}=\frac{15}{17}.
\]
每次变换的快度满足$\ee^{2\chi_1}=(1+.6)/(1-.6)=4$，故$\chi_1=\ln2$，两次同轴变换后$\chi=2\ln2=\ln4$。
\Diagram{fig:m2boostsa}{能量—动量双曲线上的快度相加}{\draw[axis](0,0)--(4.7,0)node[right]{$p/(mc)$};\draw[axis](0,0)--(0,4.8)node[above]{$E/(mc^2)$};\draw[curve,domain=0:3.8,samples=70]plot(\x,{sqrt(1+\x*\x)});\fill(.75,1.25)circle(1.5pt)node[left]{$\chi_1$};\fill(1.875,2.125)circle(1.5pt)node[above left]{$2\chi_1$};}
图\ref{fig:m2boostsa}采用$(p/(mc),E/(mc^2))$坐标。由双曲函数定义，
\[
 \gamma=\cosh\chi=\frac{4+1/4}{2}=\frac{17}{8},\qquad
 \gamma\beta=\sinh\chi=\frac{4-1/4}{2}=\frac{15}{8}.
\]
也可在第一次变换后的四速度上再作一次变换：$(\gamma_1,\gamma_1\beta_1)=(5/4,3/4)$，因此
\[
 \gamma=\frac54\frac54+\frac34\frac34=\frac{17}{8},\quad
 \gamma\beta=2\frac54\frac34=\frac{15}{8}.
\]
两项满足$\gamma^2-(\gamma\beta)^2=1$，商为$15/17$。快度相加以两次共线为前提；非共线速度不能把速度大小各自换成快度后直接相加。
\end{solution}

\section{从守恒定律推导火箭方程}\label{th:m2rocket}
火箭喷气是变质量系统，必须先选封闭的微小过程：燃烧前火箭，与燃烧后的火箭加刚排出的燃气。取燃烧前火箭的瞬时静止系，火箭静质量为$m$；燃烧后为$m+\dd m$，其中$\dd m<0$，获得微小速度$\dd v_*$。排出物沿后方以固定速度$u$运动，其静质量为$\dd\mu>0$。

在一阶精度下，火箭动能是$(\dd v_*)^2$阶，可在能量式中略去；燃气自身的相对论因子$\gamma_u$却不能略去。因此
\begin{align*}
 mc^2&=(m+\dd m)c^2+\gamma_u\dd\mu c^2,\\
 0&=m\dd v_*-\gamma_u u\dd\mu.
\end{align*}
第一式给$\gamma_u\dd\mu=-\dd m$，代入第二式得$m\dd v_*=-u\dd m$。火箭在相邻瞬时静止系间的快度变化为$\dd\chi=\dd v_*/c$，所以$\dd\chi=-(u/c)\dd m/m$。用初值$m=M_0,\chi=0$积分，
\[
 \int_0^\chi\dd q=-\frac uc\int_{M_0}^{m}\frac{\dd\mu}{\mu}
 =-\frac uc[\ln\mu]_{M_0}^{m}.
\]
由此得到
\begin{equation}\label{eq:m2rocket}
 \dd\chi=-\frac uc\frac{\dd m}m,\qquad
 \chi=\frac uc\ln\frac{M_0}m,\qquad
 v=c\tanh\left(\frac uc\ln\frac{M_0}m\right).
\end{equation}
式中$M_0$包含全部初始燃料和飞船。忽略外力、喷流方向共线、$u$在瞬时火箭系恒定，是此式的适用条件。把排出物静质量误写成$-\dd m$，会漏掉释放为喷流动能的质量能。

\Diagram{fig:m2rocketbalance}{瞬时静止系中，喷流和剩余火箭构成守恒系统}{\draw[curve](2,0)rectangle(4,1);\node at(3,.5){$m+\dd m$};\draw[axis](4.2,.5)--(5.5,.5)node[above]{$\dd v_*$};\fill(1,.5)circle(2pt);\draw[axis](.8,.5)--(-.4,.5)node[above]{$u$};\node[below]at(1,.2){$\dd\mu$};}

\begin{example}\label{ex:m2exhaust}
排出物在火箭系速度$u=0.8c$。若飞船静质量减少$1\,\mathrm{kg}$，求排出物的静质量、动能和总能量，并求燃烧前质量为$1000\,\mathrm{kg}$的飞船获得的近似快度增量。
\Diagram{fig:m2exhaustq}{质量能在喷流静能和动能之间分配}{\draw[curve](0,0)rectangle(2.2,1.1);\node at(1.1,.55){$-\dd m=1$ kg};\draw[->](2.4,.55)--(3.5,.55);\draw[curve](3.7,0)rectangle(6.2,1.1);\node at(4.95,.55){喷流总能};}
\end{example}
\begin{solution}
题中的$1\,\mathrm{kg}$损失相对于$1000\,\mathrm{kg}$很小。先按连续喷流的瞬时船系计算。设船的静质量改变$\dd m<0$，排出静质量$\dd\mu>0$。图\ref{fig:m2exhaustq}的一阶能量守恒为
\[
 -\dd m=\gamma_u\dd\mu,\qquad \gamma_u=\frac1{\sqrt{1-.8^2}}=\frac53.
\]
因此整个燃烧过程中排出静质量的和为$.6\,\mathrm{kg}$；把每一微元在其发射瞬间船系测得的能量相加，总能量为$1\,\mathrm{kg}\,c^2$，其中静能$.6\,\mathrm{kg}\,c^2$、动能$.4\,\mathrm{kg}\,c^2$。
\Diagram{fig:m2exhausta}{喷流能量分成静能与动能；两部分之和由质量损失决定}{\draw[curve](0,0)rectangle(5,1);\draw[curve](3,0)--(3,1);\node at(1.5,.5){$0.6c^2$ kg};\node at(4,.5){$0.4c^2$ kg};\node[above]at(1.5,1){静能};\node[above]at(4,1){动能};}
图\ref{fig:m2exhausta}的能量分割采用上述局部计量。喷流在瞬时船系中向后运动，动量大小为$\gamma_u u\dd\mu=-u\dd m$；飞船获得向前的动量$mc\dd\chi$，从而
\begin{align*}
 \dd\chi&=-\frac uc\frac{\dd m}{m},\\
 \Delta\chi&=-.8\int_{1000}^{999}\frac{\dd m}{m}
 =-.8[\ln m]_{1000}^{999}\\
 &=.8\ln\frac{1000}{999}=8.0040\times10^{-4}.
\end{align*}
若只保留损失比例的一阶，便为$8.0\times10^{-4}$。有限过程在初始实验室系的总喷流能量略有不同：末船能量为$999\,\mathrm{kg}\,c^2\cosh\Delta\chi$，所以
\[
 E_{\rm jet}^{\rm lab}
 =[1000-999\cosh\Delta\chi]\,\mathrm{kg}\,c^2
 \simeq.999680\,\mathrm{kg}\,c^2.
\]
约$.000320\,\mathrm{kg}\,c^2$成为船的反冲动能。各次发射的瞬时船系并不相同，局部能量相加后仍须区别于固定实验室系能量。喷流静质量的和为$.6\,\mathrm{kg}$。
\end{solution}

\begin{example}\label{ex:m2fuel}
比较两种火箭由静止达到$0.01c$所需的初末静质量比：一种喷气速度$3000\,\mathrm{m/s}$，另一种把排出能量全部变为向后的光。取$c=3.0\times10^8\,\mathrm{m/s}$。
\Diagram{fig:m2fuelq}{相同末速度下，喷气和光子排气的动量携带能力不同}{\draw[curve](0,0)rectangle(2,1);\draw[axis](-.1,.5)--(-1.3,.5)node[above]{$u$};\draw[curve](4,0)rectangle(6,1);\draw[decorate,decoration={snake,amplitude=1.2pt,segment length=5pt}](3.9,.5)--(2.7,.5);\node[above]at(3.3,.7){光};}
\end{example}
\begin{solution}
所需快度$\chi=\operatorname{artanh}(.01)=\tfrac12\ln(1.01/.99)\simeq0.01000033$。图\ref{fig:m2fuelq}左侧喷气$u/c=10^{-5}$，所以
\[
 \ln\frac{M_0}{M_f}=\frac c u\chi\simeq1000.033,
 \quad \log_{10}\frac{M_0}{M_f}\simeq434.31.
\]
质量比约为$2.0\times10^{434}$，这个模型中的困难来自低喷气速度。不能把$0.01c$称为足够低速就忽略连续变质量过程中质量比的指数。
\Diagram{fig:m2fuela}{光子火箭初末质量和快度的关系}{\draw[axis](0,0)--(4.7,0)node[right]{$\chi$};\draw[axis](0,0)--(0,3.4)node[above]{$M_0/M_f$};\draw[curve,domain=0:1.05,samples=50]plot({3*\x},{exp(\x)});\node[left]at(0,1){$1$};\node[right]at(3.1,2.8){$\ee^\chi$};}
光子火箭的排出物没有静质量，应直接以能量$\dd E_\gamma=-c^2\dd m$、动量$\dd E_\gamma/c$写守恒，得到式\eqref{eq:m2rocket}的$u=c$极限。图\ref{fig:m2fuela}给
\[
 \frac{M_0}{M_f}=\ee^\chi=\sqrt{\frac{1+\beta}{1-\beta}}\simeq1.01005.
\]
这是假定完全定向、完全质量能转化的理想模型，不包含推进装置的质量和效率限制。
\end{solution}

\begin{example}\label{ex:m2burn}
理想光子火箭从实验室原点由静止出发，固有质量损失率恒定为$-\dd m/\dd\tau=a>0$。求$m(\tau)$、$v(\tau)$、世界线和固有加速度。初质量为$M_0$。
\Diagram{fig:m2burnq}{质量在固有时间内线性减少}{\draw[axis](0,0)--(4.5,0)node[right]{$\tau$};\draw[axis](0,0)--(0,3)node[above]{$m$};\draw[curve](0,2.5)--(3.7,0);\node[left]at(0,2.5){$M_0$};\node[below]at(3.7,0){$M_0/a$};}
\end{example}
\begin{solution}
图\ref{fig:m2burnq}中$m=M_0-a\tau$，定义$s=m/M_0=1-a\tau/M_0$。物理区间为$0<s\leq1$，实际还须在干质量处停机。由光子火箭式，$\chi=\ln(M_0/m)=-\ln s$，所以
\[
 \gamma=\tfrac12(s^{-1}+s),\quad \gamma\beta=\tfrac12(s^{-1}-s),\quad
 \beta=\frac{1-s^2}{1+s^2}.
\]
世界线对固有时的导数为$\dd(ct)/\dd\tau=c\gamma$、$\dd x/\dd\tau=c\gamma\beta$。分别从0积分至$\tau$，利用
\[
 \int_0^\tau\frac{\dd q}{1-aq/M_0}
 =\left[-\frac{M_0}a\ln(1-aq/M_0)\right]_0^\tau
 =-\frac{M_0}a\ln s,
\]
\[
 \int_0^\tau(1-aq/M_0)\dd q
 =\left[q-\frac{aq^2}{2M_0}\right]_0^\tau
 =\tau-\frac{a\tau^2}{2M_0},
\]
得到
\begin{align*}
 ct&=\frac c2\left[-\frac{M_0}a\ln s+\tau-\frac{a\tau^2}{2M_0}\right],\\
 x&=\frac c2\left[-\frac{M_0}a\ln s-\tau+\frac{a\tau^2}{2M_0}\right].
\end{align*}
\Diagram{fig:m2burna}{世界线趋近光速；接近质量耗尽时固有加速度增大}{\draw[axis](-.2,0)--(4.8,0)node[right]{$ax/(cM_0)$};\draw[axis](0,-.2)--(0,4.8)node[above]{$at/M_0$};\draw[curve,domain=.02:1,samples=90]plot({(-ln(\x)-(1-\x*\x)/2)/2},{(-ln(\x)+(1-\x*\x)/2)/2});\draw[guide](0,.5)--(3.7,4.2);}
图\ref{fig:m2burna}的两个无量纲坐标之差为$(1-s^2)/2$，趋向有限值$1/2$，而二者各自对数发散。固有加速度
\[
 \alpha=c\frac{\dd\chi}{\dd\tau}
 =-\frac c s\frac{\dd s}{\dd\tau}=\frac{ca}{m},\qquad
 \frac{\dd v}{\dd t}=\frac{\alpha}{\gamma^3}.
\]
初始$\alpha=ca/M_0$，与喷光动量通量除以初质量吻合。若把$\ee^\chi$写成$1+a\tau/M_0$，就不满足质量线性减少的条件。
\end{solution}

\section{延迟时间与视运动}\label{th:m2retarded}
固定观察者在$\vct r_D$，物体在发光时刻$t_e$位于$\vct r(t_e)$，接收时刻满足
\begin{equation}\label{eq:m2retarded}
 t_o=t_e+\frac{|\vct r_D-\vct r(t_e)|}{c}.
\end{equation}
求照片要固定$t_o$解每个发光点的$t_e$，再把光线投影到成像面。求物体长度要固定$t_e$，两种操作不可互换。若观察距离远大于物体尺度，记朝向观察者的单位矢量为$\vct n$，则$\dd t_o=(1-\vct n\cdot\vct v/c)\dd t_e$。这里光程压缩产生的快慢属于接收过程，不改变局部运动速度。

\begin{example}\label{ex:m2super}
发光斑点以$\beta c$匀速运动，与朝向远处观察者的视线夹角为$\theta$。求视横向速度及其最大值；当$\theta=45^\circ$时，求超过$c$的条件。
\Diagram{fig:m2superq}{相邻发光点的横向位移和沿视线位移}{\draw[axis](0,0)--(5.3,0)node[right]{朝观察者};\draw[very thick,->](0,0)--(3,2);\draw[guide](3,0)--(3,2);\node[right]at(3,1){$v\dd t_e\sin\theta$};\node[below]at(1.5,0){$v\dd t_e\cos\theta$};\draw(.8,0)arc(0:33.69:.8);\node at(1.2,.35){$\theta$};}
\end{example}
\begin{solution}
图\ref{fig:m2superq}中第二次发光更靠近观察者，接收时间差为$\dd t_o=(1-\beta\cos\theta)\dd t_e$。将横向位移除以它，得到
\[
 v_{\rm app}=\frac{\beta c\sin\theta}{1-\beta\cos\theta}.
\]
在$45^\circ$时，条件$\beta/\sqrt2>1-\beta/\sqrt2$给$\beta>1/\sqrt2$。
\Diagram{fig:m2supera}{最大视速度的几何条件$\cos\theta=\beta$}{\draw[axis](0,0)--(4.3,0)node[right]{$\cos\theta$};\draw[axis](0,0)--(0,3)node[above]{$v_{\rm app}/c$};\draw[curve,domain=0:1,samples=90]plot({3.5*\x},{.8*sqrt(1-\x*\x)/(1-.8*\x)});\draw[guide](2.8,0)--(2.8,1.333);\node[below]at(2.8,0){$\beta$};\node[above]at(2.8,1.333){$\beta\gamma$};}
求导得
\[
 \frac1c\frac{\dd v_{\rm app}}{\dd\theta}
 =\frac{\beta\cos\theta(1-\beta\cos\theta)-\beta^2\sin^2\theta}
 {(1-\beta\cos\theta)^2}
 =\frac{\beta(\cos\theta-\beta)}{(1-\beta\cos\theta)^2}.
\]
极大满足$\cos\theta=\beta$，代回得$v_{\rm app,max}=\beta\gamma c$，如图\ref{fig:m2supera}。信号本身仍沿光锥传播；视速度是由两次不同位置、不同发光时刻构成的接收投影速度，不能用于发送超光速信号。
\end{solution}

\begin{example}\label{ex:m2photo}
半径$r$的圆截面在自身静止系位于$x'y'$平面，以速度$V$沿$x$运动。远处观察者位于$(0,-b,0)$，$b\gg r$。取快门接收时刻$t_o=b/c$。用参数$\varphi$标记圆周，求同时截面、进入照片的发光位置和横向投影，并解释旋转角。
\Diagram{fig:m2photoq}{运动圆截面及远处沿$-y$方向观察的光线}{\draw[axis](-2,0)--(3,0)node[right]{$x$};\draw[axis](0,-2.7)--(0,2.1)node[above]{$y$};\draw[curve](0,0)ellipse(1.2 and 1.5);\draw[axis](1.5,1.2)--(2.8,1.2)node[above]{$V$};\fill(0,-2.2)circle(2pt)node[right]{$D,\ b\gg r$};\draw[guide](-1,-.8)--(0,-2.2);\draw[guide](1,.8)--(0,-2.2);}
\end{example}
\begin{solution}
圆周在$S'$中是$x'=r\cos\varphi,y'=r\sin\varphi$。在实验室固定$t$，尺缩给$x=Vt+(r/\gamma)\cos\varphi$、$y=r\sin\varphi$，所以同时截面为椭圆。图\ref{fig:m2photoq}中的观察者距发光点约为$b+y$，由式\eqref{eq:m2retarded}及固定$t_o=b/c$得$t_e=-y/c$，从而
\[
 x_e=\frac r\gamma\cos\varphi-\beta r\sin\varphi,
 \qquad y_e=r\sin\varphi.
\]
定义$\cos\alpha=1/\gamma$、$\sin\alpha=\beta$，横向投影为
\[
 x_{\rm photo}=x_e=r\cos(\varphi+\alpha),\qquad \alpha=\arcsin\beta.
\]
\Diagram{fig:m2photoa}{照片横向坐标是两项投影之和；发光点并不同时}{\draw[axis](-2.4,0)--(2.9,0)node[right]{$x_{\rm photo}$};\draw[axis](0,-1.8)--(0,2.2)node[above]{$y$};\draw[curve,domain=0:360,samples=100]plot({1.5*(.8*cos(\x)-.6*sin(\x))},{1.5*sin(\x)});\draw[guide](0,0)--(1.2,0);\node[below]at(1.2,0){$\varphi=0$};\fill(-.9,1.5)circle(1.4pt)node[above left]{$\varphi=\pi/2$};\node[right]at(1.8,1.4){发光位置集合};}
图\ref{fig:m2photoa}画的是各个延迟发光点在实验室的坐标集合，不是照片上的二维圆形。沿$-y$观察时，$y$是视线深度，成像面只保留$x$和$z$。对有标记的圆柱侧面，$x_{\rm photo}$的参数平移等效于把圆柱旋转$\alpha$后投影；它不表示物体真的转动，也不表示任何有限距离、任意形状物体的照片都恰为刚体旋转。忽略的距离二阶项为物体尺度平方除以$b$。
\end{solution}
\section{连续转向与加速参考系}
\subsection{瞬时惯性系的适用范围}
光滑类时世界线在每个事件都有瞬时静止惯性系，可用来定义固有加速度和局部推进功率。若要由这些局部惯性系建立加速坐标，还须指定空间坐标和同步方式。双生子题中赋给远处地球的同时读数可随转向而跃变，旅行者的本地钟却连续走时。

对任意有限加速过程，旅行者固有时始终为$\int\dd t\sqrt{1-v^2/c^2}$，只需已知实验室轨迹。瞬时转向模型是把加速段的持续时间和路程同时趋零的极限，保留两段惯性速度的差。若加速持续时间不可略，去返总时间、最远距离以及固有时都应包括该段积分。

在一维运动中，$\chi=\operatorname{artanh}(v/c)$。链式法则给$\dd\chi/\dd t=\dot v/[c(1-v^2/c^2)]=\gamma_v^2\dot v/c$；再乘$\dd t/\dd\tau=\gamma_v$，得到有向固有加速度$a_*=c\dd\chi/\dd\tau=\gamma_v^3\dot v$，其大小为$\alpha=|a_*|$。制动时$a_*$可为负，而$\alpha$仍非负。随着$\gamma_v$增大，有限的固有加速度可对应越来越小的实验室加速度；题设须说明控制的是哪一个量。

\subsection{恒固有加速度的坐标网}
令参考飞船的固有加速度为$\alpha>0$，设$L_\alpha=c^2/\alpha$。为与飞船保持固定局部距离，可引入坐标$(\tau,\xi)$：
\begin{align*}
 ct&=(L_\alpha+\xi)\sinh(\alpha\tau/c),\\
 x&=(L_\alpha+\xi)\cosh(\alpha\tau/c)-L_\alpha.
\end{align*}
参考飞船在$\xi=0$，此处$\tau$正是其固有时。记$\chi=\alpha\tau/c$，乘积法则给
\begin{align*}
 \dd(ct)&=\sinh\chi\dd\xi
 +(L_\alpha+\xi)\cosh\chi\frac\alpha c\dd\tau,\\
 \dd x&=\cosh\chi\dd\xi
 +(L_\alpha+\xi)\sinh\chi\frac\alpha c\dd\tau.
\end{align*}
相减平方时，混合项相等而抵消；$\dd\xi^2$的系数为$\sinh^2\chi-\cosh^2\chi=-1$，$\dd\tau^2$的系数为$(L_\alpha+\xi)^2\alpha^2/c^2$，得到
\[
 \dd s^2=c^2(1+\alpha\xi/c^2)^2\dd\tau^2-\dd\xi^2.
\]
该坐标网只取$\xi>-L_\alpha$以保持同一时间方向。固定$\xi$的其他钟在同一坐标时间内经历$\dd\tau_\xi=(1+\alpha\xi/c^2)\dd\tau$，其固有加速度为$\alpha_\xi=\alpha/(1+\alpha\xi/c^2)$。要让不同位置的点保持这种固定距离，它们不能全取相同固有加速度。

\Diagram{fig:m2rindler}{固定加速坐标的世界线是共用光状渐近线的双曲线}{\draw[axis](-1.5,0)--(4.3,0)node[right]{$x/L_\alpha$};\draw[axis](0,0)--(0,4.5)node[above]{$ct/L_\alpha$};\foreach\a in{.7,1,1.5}{\draw[curve,domain=0:1.6,samples=70]plot({\a*cosh(\x)-1},{\a*sinh(\x)});}\draw[guide](-1,0)--(3.2,4.2);\node[below]at(-1,0){$-1$};}
图\ref{fig:m2rindler}中越靠前的固定坐标点，所需固有加速度越小。若一根杆的各处在实验室同一时刻开始按完全相同速度程序加速，实验室端点距离可保持不变，但随共同速度增大，其瞬时静止系端点距离会增大。这与保持局部固有距离不变是不同的加速安排，物体内部会承受应力，不能假设任意加速都保留一个绝对刚性形状。

\subsection{加速系光信号的有限传播}
沿$\xi$方向的光满足$\dd s^2=0$，所以坐标速度为$\dd\xi/\dd\tau=\pm c(1+\alpha\xi/c^2)$。局部观察者用自身钟和局部尺测得的光速仍为$c$。取$-L_\alpha<\xi_1<\xi_2$及正号，分离变量：
\[
 \Delta\tau=\frac1c\int_{\xi_1}^{\xi_2}
 \frac{\dd\xi}{1+\alpha\xi/c^2}.
\]
令$z=1+\alpha\xi/c^2$，则$\dd\xi=(c^2/\alpha)\dd z$，上下限为$z_i=1+\alpha\xi_i/c^2>0$，故
\[
 \Delta\tau=\frac c\alpha\int_{z_1}^{z_2}\frac{\dd z}z
 =\frac c\alpha[\ln z]_{z_1}^{z_2}
 =\frac c\alpha\ln\frac{1+\alpha\xi_2/c^2}{1+\alpha\xi_1/c^2}.
\]
若两站间距$\Delta\xi=\xi_2-\xi_1$很小，展开对数得到$\Delta\tau\simeq\Delta\xi/[c(1+\alpha\xi_1/c^2)]$。只有两站还位于$|\alpha\xi/c^2|\ll1$的原点附近时，才进一步近似为$\Delta\xi/c$。靠近$\xi=-L_\alpha$时坐标传播时间发散，对应图中的光状边界；它限制持续加速观察者的信号联系范围。

\section{推进方程的能量和质量含义}
\subsection{喷流速度在哪里定义}
火箭方程中的$u$是喷流在每次燃烧瞬间的飞船静止系测量值。火箭持续变速后，不同时刻喷出的两团燃气在初始实验室中的速度不同；它们无需以某一固定实验室速度向后运动。若题设给的是实验室排气速度，必须先用速度变换换成瞬时$u$，一般会得到随$m$或$\chi$变化的微分方程。

微分能量式$-\dd m=\gamma_u\dd\mu$表示飞船静质量损失中包括喷流静能和喷流动能。这里$m$是整个剩余飞船连同未用燃料的静质量，不等于某个发动机标称干质量。燃料中可转化的化学能很小，实际喷出的静质量大部分只是原料质量的搬运；光子火箭则取极端理想情况，排出物没有静质量。

在燃烧前瞬时静止系，残余飞船获得$\dd v_*$。它的反冲动能约$m\dd v_*^2/2$，相对一阶喷流能是二阶小量，所以微分推导可以忽略。有限脉冲反冲则不能忽略；例\ref{ex:m2finitepulse}保留它之后仍给相同积分质量比，原因是发射光子的实验室能量随着火箭快度连续变化，而非每一步都在同一个静止系中计量。

\subsection{质量比、推进效率与任务边界}
若真实系统只有部分失去的质量能用于定向喷流，可把静止系单位质量损失携带的后向动量记为$\eta c(-\dd m)$，$0\leq\eta\leq1$。四动量守恒给$\dd\chi=-\eta\dd m/m$。固定排气速度的理想模型有$\eta=u/c$，向后纯光束有$\eta=1$，等量前后光束有$\eta=0$。这个定义把发动机效率、废热和喷流分布压缩为一个明确动量效率，但不能仅由“多少热量释放”确定$\eta$。

对每个$\eta>0$的推进段，有$-\dd m/m=|\dd\chi|/\eta$，积分后该段的对数质量比为$\ln(m_i/m_f)=\int|\dd\chi|/\eta$。若$\eta=0$却仍排放质量，$\dd\chi=0$不能反推出质量损失，须按实际消耗另计。若巡航段不损失质量，它只增加旅行时间和距离。制动、返航各有相应推进消耗；抛弃空壳时，还须把抛弃后的剩余质量作为下一段初质量。比较燃料消耗应采用相同末态和载荷条件。

\section{延迟成像的投影几何}
照片中每一个像素对应一条到达观察者的光线。远场近似把这些光线看作平行，成像面取垂直视线的平面。设朝向观察者的单位矢量$\vct n$固定，接收时间可去掉一个共同常数写成$t_o=t_e-\vct n\cdot\vct r(t_e)/c$。固定$t_o$首先确定发射事件，随后投影
\[
 \vct R_{\rm image}=(I-\vct n\vct n^{\mathsf T})\vct r(t_e).
\]
这个投影矩阵删去沿视线深度分量，保留横向两分量。照片只保留横向投影；不同$t_e$的发射点集合也不属于某一实验室同时截面。

\Diagram{fig:m2projection}{先沿光锥选择发射事件，再向垂直视线的平面投影}{\draw[axis](0,0)--(4.5,0)node[right]{$\vct n$};\draw[axis](0,0)--(0,3)node[above]{像面横向};\fill(3,2)circle(1.5pt)node[right]{$\vct r(t_e)$};\draw[guide](3,2)--(0,2);\fill(0,2)circle(1.5pt)node[left]{投影};\draw[->](0,0)--(3,2);}
图\ref{fig:m2projection}中的横向位移除以接收时间间隔，得到视横向速度；实际位移除以发射时间间隔，得到物体的局部运动速度。具体地，$\vct v_{\rm app,\perp}=\vct v_\perp/(1-\vct n\cdot\vct v/c)$。对亚光速源，分母为正，接收顺序与发射顺序一致；视横向速度却可超过$c$，如例\ref{ex:m2super}所示。

圆柱照片中的参数转角也是投影意义上的结果。源系圆周参数$\varphi$只表示表面标记，远场横向坐标中出现$\cos(\varphi+\alpha)$，使某些原本侧面的标记进入可见投影。要判断它是否真的可见，还需考虑不透明物体的自遮挡、前后表面和发光特性；透明发光圆环与不透明圆柱不能共用同一套可见点集合。

有限观察距离时，光线不再平行，距离展开中$R_\perp^2/(2b)$项会改变延迟时刻，透视比例也随深度变化。若只保留一阶远场项，就应把“等效旋转”限定在该精度和相应物形，不把它推广为任意物体照片的普遍刚体旋转定理。

\subsection{变质量推进的局部守恒与全局能量}
火箭方程中的$m$是船体、载荷和全部剩余燃料的静质量。前文已在燃烧前的瞬时船系证明$-\dd m=\gamma_u\dd\mu$、$mc\dd\chi=\gamma_u u\dd\mu$。以下用这两个局部守恒式核对固定实验室中的能量和动量；其中$\dd\mu$是喷流静质量，$u$是其局部后向速度大小。

局部一阶守恒可积成有限的准确火箭方程。计算固定实验室中的总喷流能量时，须先把各瞬时船系的微元变换到同一实验室系：
\[
 \dd E_{\rm jet}=\gamma_v\gamma_u\dd\mu(c^2-vu),\qquad
 \dd p_{\rm jet}=\gamma_v\gamma_u\dd\mu(v-u).
\]
对船的$E=mc^2\cosh\chi$、$p=mc\sinh\chi$用乘积法则，代入$\dd m=-\gamma_u\dd\mu$及$\dd\chi=\gamma_u u\dd\mu/(mc)$，逐项得到
\begin{align*}
 \dd E&=c^2\cosh\chi\,\dd m+mc^2\sinh\chi\,\dd\chi\\
 &=-\gamma_v\gamma_u c^2\dd\mu
   +\gamma_v\gamma_u vu\dd\mu=-\dd E_{\rm jet},\\
 \dd p&=c\sinh\chi\,\dd m+mc\cosh\chi\,\dd\chi\\
 &=-\gamma_v\gamma_u v\dd\mu
   +\gamma_v\gamma_u u\dd\mu=-\dd p_{\rm jet}.
\end{align*}
这里用了$\cosh\chi=\gamma_v$、$c\sinh\chi=\gamma_vv$；两个守恒式中的喷流量与船体量均在同一实验室系计量。
\Diagram{fig:m2jetframes}{各次喷流的局部静止系不同；总能量需统一到同一参考系}{\draw[axis](0,0)--(6.1,0)node[right]{$x$};\foreach\x/\v in{.7/v_1,2.7/v_2,4.7/v_3}{\draw[curve](\x,.3)rectangle({\x+.8},.9);\draw[->]({\x+.9},.6)--({\x+1.35},.6)node[above]{$\v$};\draw[->](\x,.6)--({\x-.5},.6);}\node[below]at(3,0){同一实验室系中的连续发射位置};}
图\ref{fig:m2jetframes}只示意沿途发射，不把几个船体画成同时存在。对有限过程，先求末态船四动量，再从初态总四动量中扣除它，常比逐一积分全部喷流更省算式；两种方法必须给出同一结果。

\section{往返过程中的钟和坐标网}
\begin{example}\label{ex:m2generalturn}
旅行者去程速率$V_1$、返程速率$V_2$，在地球系距离$l$处瞬时换向。求地球和旅行者的总时间，并求换向前后同时地球时刻的跃变。
\Diagram{fig:m2generalq}{不同去返速率使世界线两段斜率不同}{\draw[axis](-.2,0)--(3.8,0)node[right]{$x$};\draw[axis](0,-.2)--(0,5.5)node[above]{$ct$};\draw[curve](0,0)--(2,2.6)--(0,4.7);\node[right]at(2,2.6){$B$};\node[below]at(2,0){$l$};}
\end{example}
\begin{solution}
在图\ref{fig:m2generalq}中取出发事件为$A:(0,0)$，转向为$B:(t_B,l)$。两段均匀速，故
\[
 t_B=\frac l{V_1},\quad T=\frac l{V_1}+\frac l{V_2},\quad
 \tau=\frac l{\gamma_1V_1}+\frac l{\gamma_2V_2}.
\]
瞬时转向不积累有限固有时。求旅行者在$B$赋给地球的同时刻时，须分别采用转向前、后的惯性系。
\Diagram{fig:m2generala}{两段同时线在转向事件处改变倾斜方向}{\draw[axis](-.2,0)--(3.7,0);\draw[axis](0,0)--(0,5.6)node[above]{$ct$};\draw[curve](0,0)--(2,2.6)--(0,4.7);\draw[guide](0,1.061538)--(3,3.369231);\draw[guide](0,4.504762)--(3,1.647619);\node[left]at(0,1.061538){$D_-$};\node[left]at(0,4.504762){$D_+$};}
去程系中与$B$同时的事件满足$\Delta t-(V_1/c^2)\Delta x=0$。令地球事件为$(t_-,0)$，则
\[
 t_--t_B=\frac{V_1}{c^2}(0-l),\qquad
 t_-=t_B-\frac{V_1l}{c^2}.
\]
返程速度为$-V_2$，同理$t_+=t_B+V_2l/c^2$。所以图\ref{fig:m2generala}中$D_+$与$D_-$的时间差为$(V_1+V_2)l/c^2$。可把旅行者赋予地球的时间分成三部分：
\[
 \frac{l}{\gamma_1^2V_1},\qquad
 \frac{(V_1+V_2)l}{c^2},\qquad
 \frac{l}{\gamma_2^2V_2}.
\]
第一项是$A$到$D_-$，末项是$D_+$到重逢；它们与中间的同时性改变量相加，利用$1/\gamma_i^2=1-V_i^2/c^2$恰恢复$T$。地球钟没有跳走，变化的是旅行者用不同同时线选取的远处事件。
\end{solution}

\begin{example}\label{ex:m2outbackopt}
规定地球往返总时间为$T$，旅行者在中途换向一次，并全程以相同速率$v$运动。求最远距离和旅行者时间；若把时间改为两段不同速率，在相同端点条件下哪类路径能积累最多固有时？
\Diagram{fig:m2optq}{固定重逢端点、一次换向的对称路径}{\draw[axis](-.2,0)--(4,0)node[right]{$x$};\draw[axis](0,0)--(0,5.4)node[above]{$ct$};\draw[curve](0,0)--(2,2.5)--(0,5);\draw[guide](0,2.5)--(2,2.5);\node[left]at(0,5){$cT$};}
\end{example}
\begin{solution}
图\ref{fig:m2optq}对称点发生在地球时间$T/2$，所以$l=vT/2$，旅行者时间$\tau=T\sqrt{1-v^2/c^2}$。若只固定出发和重逢事件，完全不离开地球的惯性直线能积累$T$，是最大值。
\Diagram{fig:m2opta}{固定端点的世界线越偏离直线，其固有时越小}{\draw[axis](-.3,0)--(3.7,0);\draw[axis](0,0)--(0,5.4)node[above]{$ct$};\draw[guide](0,0)--(1,2.5)--(0,5);\draw[curve](0,0)--(2,2.5)--(0,5);\draw[very thick](0,0)--(0,5);}
若再规定必须经过远处某个事件$B$，则$A$到$B$、$B$到$C$各自的匀速直线分别使固有时最大，总和也最大；这由每段的积分不等式逐段得到。图\ref{fig:m2opta}中的宽折线和窄折线若转向事件不同，固有时比较应保留相同端点和途经事件的约束；它取决于整条世界线。
\end{solution}

\section{火箭的有限点火与燃料分配}
\begin{example}\label{ex:m2finitepulse}
初始静止、质量$M$的理想飞船一次向后发出能量为$E_\gamma$的光脉冲。$E_\gamma$按初始飞船系测量。发射后飞船静质量$m$。求脉冲能量、末速度，并与连续光子火箭质量比比较。
\Diagram{fig:m2pulseq}{一次有限脉冲过程须保留剩余飞船的反冲能}{\draw[curve](2,0)rectangle(4,1);\node at(3,.5){初始$M$};\draw[->](1.7,.5)--(.1,.5)node[above]{$E_\gamma$};\draw[->](4.2,.5)--(5.5,.5)node[above]{$v$};}
\end{example}
\begin{solution}
在图\ref{fig:m2pulseq}的初始船系中，向后光脉冲的动量为$-E_\gamma/c$。发射后飞船向前运动，守恒式为
\[
 Mc^2=\gamma_vmc^2+E_\gamma,\qquad
 \gamma_vmv=E_\gamma/c.
\]
用前式表示末船能量，用后式表示其动量，再作质壳平方：
\begin{align*}
 m^2c^4&=(Mc^2-E_\gamma)^2-E_\gamma^2\\
 &=M^2c^4-2Mc^2E_\gamma.
\end{align*}
故$E_\gamma=(M^2-m^2)c^2/(2M)$，末船能量为$(M^2+m^2)c^2/(2M)$。
\Diagram{fig:m2pulsea}{初始系能量分为脉冲和反冲后飞船能量}{\draw[curve](0,0)rectangle(5,1);\draw[curve](2,0)--(2,1);\node at(1,.5){$E_\gamma$};\node at(3.5,.5){$\gamma_vmc^2$};\node[above]at(2.5,1){$Mc^2$};}
图\ref{fig:m2pulsea}中，末船速率和反冲动能分别为
\[
 \beta=\frac{pc}{E}=\frac{M^2-m^2}{M^2+m^2},\qquad
 K=\frac{(M-m)^2}{2M}c^2.
\]
直接相加$E_\gamma+K=(M-m)c^2$，质量损失的能量去向因而完整闭合。再由速度式计算
\[
 \frac{1+\beta}{1-\beta}=\frac{M^2}{m^2},\qquad
 \chi=\tfrac12\ln\frac{1+\beta}{1-\beta}=\ln\frac Mm.
\]
它与连续定向光子火箭的质量比相同，因为所有排出光子的四动量共线，总和仍为类光。若向多个方向发光，辐射总四动量通常不再类光，这一有限脉冲关系就需要重算。
\end{solution}

\begin{example}\label{ex:m2twosided}
飞船在瞬时静止系发出总光功率$P$，其中向后比例为$f$、向前比例为$1-f$。保持$f$不变，求质量与快度的关系。若各向同性发光，飞船能否加速？
\Diagram{fig:m2twosidedq}{相反方向光功率的差产生净推力}{\draw[curve](2,0)rectangle(4,1);\draw[->](1.8,.5)--(.1,.5)node[above]{$fP$};\draw[->](4.2,.5)--(5.9,.5)node[above]{$(1-f)P$};}
\end{example}
\begin{solution}
图\ref{fig:m2twosidedq}总能量损失给$\dd m/\dd\tau=-P/c^2$。向后减向前的光动量率为$(2f-1)P/c$，所以瞬时推力$F_0=(2f-1)P/c$，快度率$\dd\chi/\dd\tau=F_0/(mc)$。
\Diagram{fig:m2twosideda}{同样质量损失，只有光束的不对称部分产生快度变化}{\draw[axis](0,0)--(4.3,0)node[right]{$f$};\draw[axis](0,-2)--(0,2.2)node[above]{$\chi/\ln(M_0/m)$};\draw[curve](0,-1.6)--(3.6,1.6);\node[below]at(1.8,0){$1/2$};\node[below]at(3.6,0){$1$};}
消去$\tau$后$\dd\chi=-(2f-1)\dd m/m$。以初始静止、质量$M_0$为初值，
\[
 \int_0^\chi\dd q=-(2f-1)\int_{M_0}^{m}\frac{\dd\mu}{\mu}
 =-(2f-1)[\ln\mu]_{M_0}^{m},
 \quad \chi=(2f-1)\ln\frac{M_0}{m}.
\]
这给出图\ref{fig:m2twosideda}的直线。$f=1/2$时两束相抵，质量下降而速度不变；各向同性辐射也没有静止系净动量。$f<1/2$则向后加速，符号由喷光方向确定。
\end{solution}

\begin{example}\label{ex:m2varyexhaust}
某理想火箭排气速度随剩余质量变化，$u(m)=u_0m/M_0$，$u_0<c$。由初始静止加速至剩余质量$M_f$，求末快度。与把排气速度粗略取末值的结果比较。
\Diagram{fig:m2varyq}{排气速度随质量减少而下降}{\draw[axis](0,0)--(4.3,0)node[right]{$m/M_0$};\draw[axis](0,0)--(0,2.8)node[above]{$u/u_0$};\draw[curve](0,0)--(3.5,2.3);\node[below]at(3.5,0){$1$};\node[left]at(0,2.3){$1$};}
\end{example}
\begin{solution}
式\eqref{eq:m2rocket}仍可逐微元使用。图\ref{fig:m2varyq}中的$u$随质量变化，积分为
\[
 \chi_f=\int_{M_f}^{M_0}\frac{u(m)}c\frac{\dd m}m
 =\frac{u_0}{cM_0}\int_{M_f}^{M_0}\dd m
 =\frac{u_0}{cM_0}[m]_{M_f}^{M_0}
 =\frac{u_0}c\left(1-\frac{M_f}{M_0}\right).
\]
因此$\beta_f=\tanh\chi_f$。即使把$M_f$趋向0，此控制模型的快度仍只趋$u_0/c$，不发散。
\Diagram{fig:m2varya}{真实快度是对数质量坐标下排气速度的积分}{\draw[axis](0,0)--(3.4,0)node[right]{$\ln(M_0/m)$};\draw[axis](0,0)--(0,2.8)node[above]{$u/c$};\draw[curve,domain=0:2,samples=60]plot(\x,{2.3*exp(-\x)});\draw[guide](0,{2.3*exp(-2)})--(2,{2.3*exp(-2)});\draw[guide](2,0)--(2,{2.3*exp(-2)});\node[above right]at(2,{2.3*exp(-2)}){末值近似};}
若固定取末速度$u_f=u_0M_f/M_0$，会得到$\chi_{\rm app}=(u_0/c)(M_f/M_0)\ln(M_0/M_f)$。图\ref{fig:m2varya}中末值矩形低于整条曲线，所以该近似低估快度；直接用算术平均速度也无依据，正确权重是$\dd m/m$。
\end{solution}

\begin{example}\label{ex:m2fueloptimal}
火箭的全部点火满足同一固定排气速度$u$。需先加速到巡航速度，再减速到静止，初末质量比$R=M_0/M_f$已知。求可达巡航速度；若两次点火各使用相等的绝对燃料质量，是否能恰好停止？
\Diagram{fig:m2optimalq}{两次相反点火，中间质量为$M_1$}{\draw[curve](0,0)rectangle(1.5,1);\node at(.75,.5){$M_0$};\draw[->](1.7,.5)--(2.5,.5);\draw[curve](2.7,0)rectangle(4.2,1);\node at(3.45,.5){$M_1$};\draw[->](4.4,.5)--(5.2,.5);\node[right]at(5.3,.5){$M_f$};}
\end{example}
\begin{solution}
图\ref{fig:m2optimalq}中令两次燃烧之间的质量为$M_1$。加速和制动方向相反，要求的快度变化大小相同：
\[
 \chi=\frac uc\ln\frac{M_0}{M_1}
 =\frac uc\ln\frac{M_1}{M_f}.
\]
两式相等给$M_1^2=M_0M_f$；相加给$2\chi=(u/c)\ln R$。所以$M_1=\sqrt{M_0M_f}$，巡航速度为$c\tanh[(u/2c)\ln R]$。
\Diagram{fig:m2optimala}{两次相等的是质量比；后段在更小质量下产生相同快度变化}{\draw[axis](0,0)--(5.2,0)node[right]{过程};\draw[axis](0,0)--(0,3)node[above]{$\ln m$};\draw[curve](0,2.6)--(2.3,1.5)--(4.6,.4);\draw[guide](0,1.5)--(2.3,1.5);\node[right]at(2.3,1.5){$\ln\sqrt{M_0M_f}$};}
图\ref{fig:m2optimala}按$\ln m$标出两段等长。若改为各消耗相同绝对燃料质量，则$M_1=(M_0+M_f)/2$，其最终快度为
\begin{align*}
 \chi_f&=\frac uc\left(\ln\frac{M_0}{M_1}-\ln\frac{M_1}{M_f}\right)\\
 &=\frac uc\ln\frac{4M_0M_f}{(M_0+M_f)^2}<0
 \quad(M_0>M_f).
\end{align*}
后段制动超过所需，船最终反向运动。由于每消耗一单位质量的快度增量与当前质量成反比，同样千克数的燃料在后段作用更大。满足停车条件时，两段消耗量为$M_0-\sqrt{M_0M_f}$和$\sqrt{M_0M_f}-M_f$，前段较多，其比值为$\sqrt R$。
\end{solution}

\section{近场接收与远场视速度}
\begin{example}\label{ex:m2nearretarded}
观察者在原点，发光点沿$x$轴以$v$运动，轨迹为$(vt,b)$，$b>0$。推导接收时间对发射时间的导数，说明为什么严格近场中没有单个固定观察角。
\Diagram{fig:m2nearq}{有限冲量参数$b$的直线运动}{\draw[axis](-2.5,0)--(3,0)node[right]{$x$};\draw[axis](0,0)--(0,3)node[above]{$y$};\draw[curve](-2.2,1.5)--(2.7,1.5);\fill(1.8,1.5)circle(1.5pt)node[above]{$S$};\draw[guide](0,0)--(1.8,1.5);\node[left]at(0,1.5){$b$};}
\end{example}
\begin{solution}
图\ref{fig:m2nearq}的源距$R(t_e)=\sqrt{v^2t_e^2+b^2}$，接收时间$t_o=t_e+R/c$。求导得
\[
 \frac{\dd t_o}{\dd t_e}
 =1+\frac1{2c}(v^2t_e^2+b^2)^{-1/2}(2v^2t_e)
 =1+\frac{v^2t_e}{c\sqrt{v^2t_e^2+b^2}}.
\]
若定义从源指向观察者的方向$\vct n=-(vt_e,b)/R$，就恢复$1-\vct n\cdot\vct v/c$。该导数在$1-\beta$与$1+\beta$之间，始终为正。
\Diagram{fig:m2neara}{最近点前信号间隔压缩，经过后拉长}{\draw[axis](-2.7,0)--(3,0)node[right]{$vt_e/b$};\draw[axis](0,0)--(0,2.8)node[above]{$\dd t_o/\dd t_e$};\draw[curve,domain=-2.5:2.5,samples=80]plot(\x,{1.4*(1+.8*\x/sqrt(1+\x*\x))});\draw[guide](-2.5,1.4)--(2.5,1.4);}
图\ref{fig:m2neara}在$t_e=0$导数为1，但对应接收时刻为$b/c$。由单调性，亚光速源每个接收时刻只有一个延迟发射时刻；不能把物体“现在的位置”代入传播距离后再解释照片。远场短时间近似才允许把$\vct n$视为常量。
\end{solution}

\begin{example}\label{ex:m2timecompression}
向观察者运动的光源在固有时$\Delta\tau$内完成一次亮度变化，远场观察角为$\theta$。求接收变化持续时间；求光源同一段运动对应的视横向长度。
\Diagram{fig:m2compressq}{固有持续时间、发射时间和到达时间的三层关系}{\draw[curve](0,0)rectangle(1.6,1);\node at(.8,.5){$\Delta\tau$};\draw[->](1.8,.5)--(2.6,.5);\node at(3.1,.5){$\Delta t_e$};\draw[->](3.7,.5)--(4.5,.5);\node[right]at(4.6,.5){$\Delta t_o$};}
\end{example}
\begin{solution}
令$\vct n$从源指向远处观察者，图\ref{fig:m2compressq}中的夹角为$\vct n$与$\vct v$的夹角。源连续两次发射间隔$\Delta t_e$，第二次更靠近观察者$\vct n\cdot\vct v\Delta t_e$，故在方向近似不变的条件下
\[
 \Delta t_o=\Delta t_e-\frac{\vct n\cdot\vct v}{c}\Delta t_e
 =(1-\beta\cos\theta)\gamma\Delta\tau
 =\frac{\Delta\tau}{\delta}.
\]
这里先使用时间膨胀把源的固有时换成发射时间，再计算光程的缩短，两步各有不同来源。
\Diagram{fig:m2compressa}{横向投影长度用发射期间的真实位移，视速度用接收时间相除}{\draw[axis](0,0)--(4.2,0)node[right]{视线};\draw[->,thick](0,0)--(2.7,1.8);\draw[guide](2.7,0)--(2.7,1.8);\node[right]at(2.7,.9){$v\gamma\Delta\tau\sin\theta$};}
图\ref{fig:m2compressa}的同段横向位移为$\Delta X_\perp=v\gamma\Delta\tau\sin\theta$，因此
\[
 v_{\rm app}=\frac{\Delta X_\perp}{\Delta t_o}
 =\frac{v\sin\theta}{1-\beta\cos\theta}.
\]
角度函数与例\ref{ex:m2super}相同，其求导过程已在那里给出：最大视横速出现在$\cos\theta=\beta$，为$\gamma v$。时间压缩因子$\delta$则在$\theta=0$最大，为$\gamma(1+\beta)$；此时横向位移为零。最快亮度变化与最大视横速对应不同观察角。若近距离使视线在变化期间显著转动，应沿实际运动积分微分延迟关系。
\end{solution}

\section{有限加速装置的距离与信号}
\subsection{固定固有距离、等实验室距离与雷达距离}
本讲的加速坐标已给出$\dd s^2=c^2(1+\alpha\xi/c^2)^2\dd\tau^2-\dd\xi^2$。在同一坐标时间的局部静止切片上，沿杆的长度元是$\dd\xi$。所以前后端固定在$\xi=0,L$时，杆的固有长度保持$L$；前端固有加速度却为$\alpha/(1+\alpha L/c^2)$。这是有限装置保持固有长度的一种安排，需要不同位置的动力控制配合。

另一种安排是两个飞船从实验室同时启动，始终使用相同实验室运动程序，它们的实验室间距保持原值。这不保持加速过程中的共动距离。第三种距离是某一站按发射、接收光的往返时间定义的雷达距离，它还包含信号传播期间的加速影响。题目中的“距离保持不变”必须先说明属于哪一种。

\begin{example}加速杆上的钟与后端雷达
一根保持固有长度$L$的杆沿$x$方向加速，后端$A$的固有加速度恒为$\alpha$，前端$B$固定在加速坐标$\xi=L$。后端发出光，经前端理想反射后返回。求后端测得的往返固有时间和雷达距离，并求前端向后端持续发光时的接收频率比。
\ExpansionDiagram{fig:dep-m2-radar-q}{固定加速坐标中的两站及往返信号，图示$\alpha L/c^2=1$}{
\draw[->](0,0)--(6,0)node[right]{$\xi$};\draw[thick](1,-.4)--(1,2.5)node[above]{$A:0$};\draw[thick](5,-.4)--(5,2.5)node[above]{$B:L$};\draw[->,domain=1:5,samples=70]plot(\x,{.3+ln(1+(\x-1)/4)});\draw[->,domain=5:1,samples=70]plot(\x,{.3+2*ln(2)-ln(1+(\x-1)/4)});\node[left]at(.9,.3){发射};\node[left]at(.9,1.6863){返回};\node[left]at(.9,2.3){$\tau$};}
\end{example}
\begin{solution}
在图\ref{fig:dep-m2-radar-q}的坐标中，前行光满足$\dd\xi/\dd\tau=c(1+\alpha\xi/c^2)$。分离变量，从0积分到$L$，单程所需后端坐标时间为
\[
 t_+=\int_0^L\frac{\dd\xi}{c(1+\alpha\xi/c^2)}
 =\left[\frac c\alpha\ln(1+\alpha\xi/c^2)\right]_0^L
 =\frac c\alpha\ln(1+\alpha L/c^2).
\]
反向光的坐标速度反号，但区间次序也反向，故$t_-=t_+$。后端钟的固有时就是该坐标时间，得到
\[
 \Delta\tau_A=\frac{2c}{\alpha}\ln(1+\alpha L/c^2),\qquad
 D_A=\frac{c\Delta\tau_A}{2}
 =\frac{c^2}{\alpha}\ln(1+\alpha L/c^2).
\]
\ExpansionDiagram{fig:dep-m2-radar-a}{雷达距离与固定固有长度的比较}{
\draw[->](0,0)--(4,0)node[right]{$\alpha L/c^2$};\draw[->](0,0)--(0,3.4)node[above]{$\alpha D_A/c^2$};\draw[dashed](0,0)--(3.1,3.1);\draw[thick,domain=0:3.1,samples=80]plot(\x,{ln(1+\x)});\node[right]at(2.9,1.2){$\ln(1+x)$};}
图\ref{fig:dep-m2-radar-a}中$D_A<L$。弱加速展开为$D_A=L-\alpha L^2/(2c^2)+\cdots$，这与局部测尺得到$L$不同，是两种测量程序的差别。

前端相邻两个波峰的发射固有时间间隔为$\Delta\tau_B$。固定两站之间的坐标传播时间不随发射时刻变化，因此相邻波峰的接收坐标时间间隔等于发射坐标间隔。又$\Delta\tau_B=(1+\alpha L/c^2)\Delta\tau$、$\Delta\tau_A=\Delta\tau$，所以
\[
 \frac{\nu_A}{\nu_B}=1+\frac{\alpha L}{c^2}.
\]
这里$\nu_B$按前端本地钟定义。若只讨论后端发出再由固定前端反射回来的光，两个方向的频率因子相乘抵消，返回后端的频率等于原发射频率；单程两站比较与同站往返比较不可混用。
\end{solution}

\subsection{相同加速程序造成的有限长度效应}
两端距离很短时，瞬时长度收缩关系可以估算共动长度。有限间距下，同一瞬时惯性系选出的两个端点事件往往处于不同加速阶段，须先求同时面与两条世界线的交点，再计算其空间距离。

\begin{example}始终保持实验室间距的两飞船
两飞船初始静止，实验室间距$d>0$。从$t=0$起，两者均沿$+x$保持相同固有加速度$\alpha$，并持续加速。后船$A$的快度为$\chi$时，在其瞬时静止惯性系中，与该事件同时的前船$B$事件具有何种快度？求该系中的两船距离。
\ExpansionDiagram{fig:dep-m2-bell-q}{两条世界线在实验室等时切片上的间距恒为$d$}{
\draw[->](-.3,0)--(4,0)node[right]{$x$};\draw[->](0,0)--(0,3.4)node[above]{$ct$};\draw[thick,domain=0:1.5,samples=60]plot({cosh(\x)-1},{sinh(\x)});\draw[thick,domain=0:1.5,samples=60]plot({cosh(\x)-.4},{sinh(\x)});\draw[dashed](.543,1.175)--(1.143,1.175)node[midway,above]{$d$};\node[below]at(0,0){$A$};\node[below]at(.6,0){$B$};}
\end{example}
\begin{solution}
设$\ell_\alpha=c^2/\alpha$。图\ref{fig:dep-m2-bell-q}的两条世界线参数为
\[
 ct_A=\ell_\alpha\sinh\chi,\quad
 x_A=\ell_\alpha(\cosh\chi-1),
\]
\[
 ct_B=\ell_\alpha\sinh\eta,\quad
 x_B=d+\ell_\alpha(\cosh\eta-1).
\]
在$A$的瞬时系中同时，要求$\Delta(ct)'=\cosh\chi\,\Delta(ct)-\sinh\chi\,\Delta x=0$。将两条世界线代入，使用双曲正弦差角公式得到
\[
 \ell_\alpha\sinh(\eta-\chi)=d\sinh\chi,
 \qquad \eta=\chi+\operatorname{arsinh}
 \left(\frac d{\ell_\alpha}\sinh\chi\right).
\]
这表明前船上所选的同时事件晚于实验室等时所选事件，且在持续加速期间已经具有更大的快度。
\ExpansionDiagram{fig:dep-m2-bell-a}{后船瞬时系的同时线与前船世界线相交于更晚事件}{
\draw[->](-.3,0)--(4,0)node[right]{$x$};\draw[->](0,0)--(0,3.4)node[above]{$ct$};\draw[thick,domain=0:1.8,samples=70]plot({cosh(\x)-1},{sinh(\x)});\draw[thick,domain=0:1.8,samples=70]plot({cosh(\x)-.4},{sinh(\x)});\draw[dashed](.543,1.175)--(3.5,1.175);\draw[->](.543,1.175)--(2.9,2.97)node[right]{同时线};\fill(.543,1.175)circle(1.5pt)node[left]{$A$};\fill(2.324,2.533)circle(1.5pt)node[right]{$B$};}
距离须变换空间分量：
\begin{align*}
 D_A&=\cosh\chi\,\Delta x-\sinh\chi\,\Delta(ct)\\
 &=d\cosh\chi+\ell_\alpha[\cosh(\eta-\chi)-1]\\
 &=d\cosh\chi+\ell_\alpha\left[
 \sqrt{1+(d/\ell_\alpha)^2\sinh^2\chi}-1\right].
\end{align*}
图\ref{fig:dep-m2-bell-a}中的同时线不是实验室水平切片，所以有限间距的精确结果比$d\cosh\chi$还大。当$(d/\ell_\alpha)|\sinh\chi|\ll1$时，第二项约为$d^2\sinh^2\chi/(2\ell_\alpha)$，此时才可用$D_A\simeq d\gamma$近似。若两船在某实验室时刻停止加速，则前船上所取事件可能已经进入匀速段，必须分段接续世界线，不能继续沿用此式。
\end{solution}

\section{同一实验室时刻怎样读取加速杆两端的钟}
比较加速杆两端的钟，需要先指定读取事件。加速坐标的同时面和实验室的同时面选取不同的事件对，钟面读数也随之改变。以下令杆的固有长度保持不变，比较实验室在同一时刻读取的两端固有时。

设杆后端固有加速度为$\alpha$，固有长度$L$，记$s=1+\alpha L/c^2>1$。本讲加速坐标给后端$ct=(c^2/\alpha)\sinh(\alpha\tau/c)$；前端则为$ct=s(c^2/\alpha)\sinh(\alpha\tau/c)$。在两端都取实验室$t=0$读数为零后，前端固有时还要乘局部钟率因子$s$。

\begin{example}同时拍摄两只加速钟的坐标读数
理想记录装置已经消除光传播延迟，得到实验室同一时刻$t$的后端钟读数$T_A$和前端钟读数$T_B$。两钟按上述保持固有长度的方案从静止开始持续加速。求两读数，并判断哪只钟读数更大。这里讨论坐标同时读数，不讨论未经延迟修正的照片。
\ExpansionDiagram{fig:dep-m2-clock-q}{实验室水平同时线选取两条加速世界线上的不同加速坐标时刻}{
\draw[->](0,0)--(4,0)node[right]{$x$};\draw[->](0,0)--(0,3.8)node[above]{$ct$};\draw[thick,domain=0:1.8,samples=70]plot({cosh(\x)-1},{sinh(\x)});\draw[thick,domain=0:1.6,samples=70]plot({1.5*cosh(\x)-1},{1.5*sinh(\x)});\draw[dashed](0,2)--(3.3,2)node[right]{同一$t$};\node[below]at(0,0){$A$};\node[below]at(.5,0){$B$};}
\end{example}
\begin{solution}
图\ref{fig:dep-m2-clock-q}要求先分别反解世界线时间。记$s=1+\alpha L/c^2>1$。后端固有时就是其加速坐标时间，前端固有时为局部钟率因子$s$乘其坐标时间，所以
\[
 T_A=\frac c\alpha\operatorname{arsinh}\frac{\alpha t}{c},\qquad
 T_B=\frac{sc}{\alpha}\operatorname{arsinh}\frac{\alpha t}{sc}.
\]
不能把$T_B$直接写为$sT_A$，因为那样比较的是相同加速坐标时间，不是图中的实验室水平切片。
\ExpansionDiagram{fig:dep-m2-clock-a}{初始读数相同，随后前端钟在实验室相同时间内积累更多固有时}{
\draw[->](0,0)--(5,0)node[right]{$\alpha t/c$};\draw[->](0,0)--(0,3)node[above]{$\alpha T/c$};\draw[thick,domain=0:4.5,samples=90]plot(\x,{ln(\x+sqrt(1+\x^2))});\draw[dashed,domain=0:4.5,samples=90]plot(\x,{1.5*ln(\x/1.5+sqrt(1+(\x/1.5)^2))});\node[right]at(3.5,2.7){$T_B$};\node[right]at(3.5,1.7){$T_A$};}
两式求导，
\[
 \frac{\dd T_A}{\dd t}
 =\frac c\alpha\frac{\alpha/c}{\sqrt{1+(\alpha t/c)^2}}
 =\frac1{\sqrt{1+(\alpha t/c)^2}},
\]
\[
 \frac{\dd T_B}{\dd t}
 =\frac{sc}\alpha\frac{\alpha/(sc)}{\sqrt{1+(\alpha t/(sc))^2}}
 =\frac1{\sqrt{1+(\alpha t/(sc))^2}}.
\]
因$s>1$，对$t>0$有$\dot T_B>\dot T_A$，结合初始两读数为零，得到$T_B>T_A$。图\ref{fig:dep-m2-clock-a}画$s=1.5$的曲线。初期展开可得$T_B-T_A\simeq\alpha^2t^3(1-s^{-2})/(6c^2)$，差值不是从启动瞬间就线性增长。

物理上，保持固有长度要求前端加速度较小。在同一实验室时刻，前端速率较低，钟慢效应较弱，所以固有时积累更多。这一解释依赖本题的具体加速安排；若两端使用相同固有加速度，世界线和所比较距离都改变，不能照搬这些读数。
\end{solution}

\section{习题与解答}
\begin{exercise}\label{ex:m2horizon}
飞船从地球由静止出发，持续保持正向固有加速度$\alpha$。地球在$t_e>0$发出追赶飞船的光。求飞船接收该信号的固有时，并判断地球所有未来信号是否都能追上。
\Diagram{fig:m2horizonq}{恒固有加速世界线及追赶它的光线}{\draw[axis](-.2,0)--(4.5,0)node[right]{$\alpha x/c^2$};\draw[axis](0,0)--(0,4.6)node[above]{$\alpha t/c$};\draw[curve,domain=0:2,samples=90]plot({cosh(\x)-1},{sinh(\x)});\draw[->](0,.5)--(.25,.75);\draw[guide](0,1)--(3.4,4.4);\node[left]at(0,1){$1$};}
\end{exercise}
\begin{solution}
飞船世界线$ct=L_\alpha\sinh\chi$、$x=L_\alpha(\cosh\chi-1)$，$\chi=\alpha\tau/c$。图\ref{fig:m2horizonq}的信号满足$x=c(t-t_e)$，代入后
\[
 \frac{\alpha t_e}c=\sinh\chi-\cosh\chi+1=1-\ee^{-\chi}.
\]
所以仅当$0<t_e<c/\alpha$有有限接收固有时
\[
 \tau_R=-\frac c\alpha\ln(1-\alpha t_e/c).
\]
\Diagram{fig:m2horizona}{发射时刻逼近$c/\alpha$时，接收固有时发散}{\draw[axis](0,0)--(4.5,0)node[right]{$\alpha t_e/c$};\draw[axis](0,0)--(0,3.8)node[above]{$\alpha\tau_R/c$};\draw[curve,domain=0:.96,samples=90]plot({3.5*\x},{-ln(1-\x)});\draw[guide](3.5,0)--(3.5,3.5);\node[below]at(3.5,0){$1$};}
图\ref{fig:m2horizona}在边界发散。$t_e\geq c/\alpha$的信号永不追上持续加速的飞船，虽然飞船在每个有限时刻速度都小于$c$。这来自其未来世界线不断趋近一条光状渐近线。若飞船后来停机或转向，未来轨迹改变，不能继续使用永久加速的信号边界。
\end{solution}

\begin{exercise}\label{ex:m2labconstant}
光子火箭从静止出发，按初始实验室时间计算的后向光功率恒为$P$。求火箭剩余静质量和速度随实验室时间的变化，并与恒固有功率模型区分。初始质量$M_0$。
\Diagram{fig:m2labconstantq}{功率固定在同一实验室中定义}{\draw[curve](2,0)rectangle(4,1);\draw[->](1.8,.5)--(.1,.5)node[above]{$P$实验室};\draw[->](4.2,.5)--(5.7,.5)node[above]{$v(t)$};}
\end{exercise}
\begin{solution}
图\ref{fig:m2labconstantq}截至时间$t$发出的总光能$Pt$、总后向动量$-Pt/c$。全系统守恒给飞船总能量$E=M_0c^2-Pt$、动量$p=Pt/c$。令$s=Pt/(M_0c^2)$，质壳关系得到
\[
 m=M_0\sqrt{1-2s},\qquad \beta=\frac{pc}E=\frac s{1-s},\quad 0\leq s<1/2.
\]
实际还须在$m$达到干质量时停机，不能把$m=0$极限当有质量飞船状态。
\Diagram{fig:m2labconstanta}{实验室定功率与固有定功率对应不同质量时间关系}{\draw[axis](0,0)--(4.4,0)node[right]{$s$};\draw[axis](0,0)--(0,2.8)node[above]{$m/M_0$};\draw[curve,domain=0:.499,samples=90]plot({7*\x},{2.3*sqrt(1-2*\x)});\node[below]at(3.5,0){$1/2$};\node[left]at(0,2.3){$1$};}
图\ref{fig:m2labconstanta}的质量不是实验室时间的线性函数。由$P=P_*(1-\beta)$，所需固有功率$P_*=P(1-s)/(1-2s)$会增大。固有时间满足$\dd\tau=\sqrt{1-2s}\dd t/(1-s)$。令$K=M_0c^2/P$、$w=\sqrt{1-2s}$，则$\dd t=K\dd s=-Kw\dd w$、$1-s=(1+w^2)/2$，所以
\[
 \tau=-2K\int_1^w\frac{q^2}{1+q^2}\dd q
 =-2K[q-\arctan q]_1^w.
\]
代入上下限，得到
\[
 \tau=\frac{2M_0c^2}P\left[1-w+\arctan w-\frac\pi4\right].
\]
在$t=0$有$w=1,\tau=0$。可见恒功率条件随参考系而变。固有功率恒定时采用$m=M_0-P_*\tau/c^2$，对应另一种推进程序。
\end{solution}

\begin{exercise}\label{ex:m2midburn}
光子火箭需从静止到达距离$L$处并停止，前半程保持固有加速度$\alpha$，后半程保持同样大小的固有减速度。求地球时间、船上时间和初末质量比。
\Diagram{fig:m2midq}{中点反转推力，速度连续}{\draw[axis](-.3,0)--(5.8,0)node[right]{$x$};\fill(0,0)circle(1.5pt);\fill(5,0)circle(1.5pt);\draw[guide](2.5,-.7)--(2.5,1.3);\draw[->](.5,.6)--(1.8,.6)node[above]{$\alpha$};\draw[->](4.5,.6)--(3.2,.6)node[above]{$-\alpha$};}
\end{exercise}
\begin{solution}
对图\ref{fig:m2midq}的前半程，用恒固有加速度世界线，令$q=\alpha L/(2c^2)$，得到中点快度$\chi_m=\operatorname{arcosh}(1+q)$。该段地球时间$c\sinh\chi_m/\alpha$、船上时间$c\chi_m/\alpha$。
\Diagram{fig:m2mida}{固有时中的快度先线性增加再线性减少}{\draw[axis](0,0)--(5.2,0)node[right]{$\tau$};\draw[axis](0,0)--(0,3)node[above]{$\chi$};\draw[curve](0,0)--(2.3,2.5)--(4.6,0);\node[above]at(2.3,2.5){$\chi_m$};}
后半程是时间反向的减速过程，图\ref{fig:m2mida}两段持续时间相同。因此$T=2c\sqrt{q(2+q)}/\alpha$、$\tau=2c\operatorname{arcosh}(1+q)/\alpha$。光子火箭消耗与总绝对快度$2\chi_m$对应，$M_0/M_f=\ee^{2\chi_m}=[1+q+\sqrt{q(2+q)}]^2$。减速阶段仍耗质量，不能让快度符号变化抵消燃料消耗。
\end{solution}

\begin{exercise}\label{ex:m2staging}
固定排气速度$u$的火箭初质量100，燃烧至60时丢弃质量20的空壳，再燃烧至最终质量20。求总快度，并与同样初质量、最终携带40质量但不丢壳的单级模型比较。所有质量使用同一单位。
\Diagram{fig:m2stageq}{无相对弹射的丢壳瞬间，飞船速度不变}{\node at(0,0){$100$};\draw[->](.6,0)--(1.5,0);\node at(2,0){$60$};\draw[->](2.6,0)--(3.5,0);\node at(4,0){$40$};\draw[->](4.6,0)--(5.5,0);\node at(6,0){$20$};\node[below]at(3,-.6){丢壳20};}
\end{exercise}
\begin{solution}
图\ref{fig:m2stageq}第一段快度$(u/c)\ln(100/60)$。丢壳时空壳与飞船相对静止，四动量按共同四速度分开，没有快度突变；第二段快度$(u/c)\ln(40/20)$。总快度$\chi=(u/c)\ln(10/3)$。
\Diagram{fig:m2stagea}{对数质量图中竖直丢壳不贡献推进快度}{\draw[axis](0,0)--(5.2,0)node[right]{阶段};\draw[axis](0,0)--(0,3.5)node[above]{$\ln m$};\draw[curve](0,3)--(1.8,2.3)--(1.8,1.7)--(4.2,.7);\node[right]at(1.8,2){丢壳};}
不丢壳的对照模型结束时含壳和有效载荷总质量40。其快度为
\[
 \chi_1=\frac uc\ln\frac{100}{40},\qquad
 \chi-\chi_1=\frac uc\ln\frac43>0.
\]
图\ref{fig:m2stagea}对应的分级增益为后一式。若把无推进丢掉的20也按火箭式积分，会错误地多算一段快度；比较模型的有效载荷、燃料和空壳质量必须逐项一致。
\end{solution}

\begin{exercise}\label{ex:m2labpower}
光子火箭沿$+x$运动，瞬时静止系向后辐射功率为$P_*$。求按实验室发射时间计算的光功率，并由它和飞船动能变化核对质量能守恒。
\Diagram{fig:m2labpowerq}{后向光子在实验室红移}{\draw[curve](2,0)rectangle(4,1);\draw[axis](4.2,.5)--(5.6,.5)node[above]{$V$};\draw[axis](1.8,.5)--(.2,.5)node[above]{$P_*$};}
\end{exercise}
\begin{solution}
图\ref{fig:m2labpowerq}后向光子能量因子为$\gamma(1-\beta)$，发射时间因子为$\gamma$，所以$P_\gamma=P_*(1-\beta)$。同时$\dd m/\dd t=-P_*/(\gamma c^2)$，飞船快度率$\dd\chi/\dd t=P_*/(\gamma mc^2)$。
\Diagram{fig:m2labpowera}{飞船总能量减小率等于后向辐射实验室功率}{\draw[curve](0,0)rectangle(2.2,1);\node at(1.1,.5){$\gamma mc^2$};\draw[->](2.4,.5)--(3.5,.5);\node[right]at(3.7,.5){$P_*(1-\beta)$};}
对图\ref{fig:m2labpowera}的飞船总能量求导，利用$\dd\gamma/\dd\chi=\gamma\beta$，有
\[
 \frac{\dd(\gamma mc^2)}{\dd t}
 =c^2\left(\gamma\frac{\dd m}{\dd t}
 +m\frac{\dd\gamma}{\dd\chi}\frac{\dd\chi}{\dd t}\right)
 =-P_*+mc^2\gamma\beta\frac{P_*}{\gamma mc^2}
 =-P_*(1-\beta).
\]
飞船虽加速，总能量仍因辐射流失而下降；静质量损失的一部分转为更大的运动因子。不能把“速度增加”直接等同于飞船总能量必增。
\end{solution}

\begin{exercise}\label{ex:m2relay}
沿用地球与恒星相距10光年、速度$0.8c$的过程，但旅行者不转向。另一只返航钟在恒星处与其相遇，把读数设置为7.5年后飞回地球。比较重逢读数与原双生子结果，解释是否必须由同一只钟经历加速才能得到差值。
\Diagram{fig:m2relayq}{两只匀速钟在恒星处交接读数}{\draw[axis](-.2,0)--(4.2,0);\draw[axis](0,0)--(0,5.4)node[above]{$ct$};\draw[curve](0,0)--(2,2.5)--(3.7,4.625);\draw[curve](3.6,.5)--(0,5);\fill(2,2.5)circle(1.4pt)node[right]{$B$交接};}
\end{exercise}
\begin{solution}
图\ref{fig:m2relayq}去程钟从出发到$B$经历7.5年，返航钟从$B$到地球也经历7.5年。它在$B$被设置为7.5年，所以到地球读15年，地球仍为25年。
\Diagram{fig:m2relaya}{交接后的读数拼接两段不同世界线的固有时间}{\draw[axis](0,0)--(5.1,0)node[right]{地球时间};\draw[axis](0,0)--(0,3)node[above]{拼接读数};\draw[curve](0,0)--(4.6,2.4);\draw[guide](2.3,0)--(2.3,1.2);\node[below]at(2.3,0){12.5年};\node[right]at(4.6,2.4){15年};}
图\ref{fig:m2relaya}的拼接读数沿两段组成的折线累积，几何上与瞬时转向世界线相同。没有单只返航钟需要在$B$改变速度，但必须实际校准读数；两只钟若早在各自惯性系随意同步，初始偏移未必满足本题。差值由选择的时空路径决定，加速是连接路径的一种实现方式。
\end{solution}

\begin{exercise}\label{ex:m2retardedzero}
发光点在实验室沿$(vt,b)$运动，观察者在原点。观察者在$t_o=0$拍照，而该坐标时刻物体恰在最近点$x=0$。求照片中光的发射时刻和发射位置。
\Diagram{fig:m2zeroq}{同时最近点与照片发光点不是同一点}{\draw[axis](-3,0)--(2.5,0)node[right]{$x$};\draw[axis](0,0)--(0,2.8)node[above]{$y$};\draw[curve](-2.6,1.6)--(2.2,1.6);\fill(0,1.6)circle(1.5pt)node[above right]{同时位置};}
\end{exercise}
\begin{solution}
图\ref{fig:m2zeroq}的延迟条件为$0=t_e+\sqrt{v^2t_e^2+b^2}/c$，要求$t_e<0$。移项平方后$c^2t_e^2=v^2t_e^2+b^2$，所以$t_e=-\gamma b/c$、$x_e=-\beta\gamma b$。
\Diagram{fig:m2zeroa}{入眼光线来自最近点之前，路程为$\gamma b$}{\draw[axis](-3.2,0)--(2.3,0)node[right]{$x$};\draw[axis](0,0)--(0,2.7)node[above]{$y$};\fill(-2,1.6)circle(1.5pt)node[above]{$x_e=-\beta\gamma b$};\draw[->](-2,1.6)--(0,0);\draw[guide](0,0)--(0,1.6);\node[right]at(0,.8){$b$};}
代回原式，光程$R_e=\sqrt{\beta^2\gamma^2b^2+b^2}=\gamma b$，恰在$\gamma b/c$后到达，验证选取了正确负根。图\ref{fig:m2zeroa}中的源虽在坐标时间0已到最近点，照片仍显示此前位置；照片记录的是经过传播后到达的旧信号。
\end{solution}

\begin{exercise}\label{ex:m2fuelbudget}
理想光子火箭总初末质量比为$R$，需把给定质量的有效载荷送到远方并相对地球停下。不规定距离，忽略巡航损失。求最大巡航洛伦兹因子以及质量比达到100时的巡航速度。
\Diagram{fig:m2budgetq}{燃料预算分成加速和减速两个相等快度段}{\draw[axis](0,0)--(5.2,0);\draw[curve](0,.3)--(2.3,1.8)--(4.6,.3);\node[above]at(2.3,1.8){巡航};\node[below]at(0,0){静止};\node[below]at(4.6,0){静止};}
\end{exercise}
\begin{solution}
图\ref{fig:m2budgetq}总绝对快度为$\ln R$，要最终停止，两段均为$\chi=\tfrac12\ln R$。所以$\gamma_{\max}=\cosh(\tfrac12\ln R)=(\sqrt R+1/\sqrt R)/2$，$\beta=(R-1)/(R+1)$。
\Diagram{fig:m2budgeta}{保留制动燃料会显著降低同一初末质量比的巡航速度}{\draw[axis](0,0)--(4.8,0)node[right]{$R$};\draw[axis](0,0)--(0,2.8)node[above]{$\beta$};\draw[curve,domain=1:100,samples=80]plot({.04*\x},{2.3*(\x-1)/(\x+1)});\node[below]at(4,0){100};}
取$R=100$，$\gamma=5.05$、$\beta=99/101\simeq.980198$，如图\ref{fig:m2budgeta}。若把全部燃料只用来加速，速度可达$(R^2-1)/(R^2+1)$，但剩余质量不足以再完成制动。任务若要求到达后停下，质量比还要包含制动段的消耗。
\end{solution}

\begin{exercise}\label{ex:m2signals}
例\ref{ex:m2twins}中地球每年发出一束时间标记光信号。求旅行者去程和返程接收到的频率，并据此还原地球25年的连续钟信号。
\Diagram{fig:m2signalsq}{地球连续发光，旅行者先远离后靠近}{\draw[axis](-.2,0)--(3.5,0)node[right]{$x$};\draw[axis](0,-.2)--(0,5.5)node[above]{$ct$};\draw[curve](0,0)--(2,2.5)--(0,5);\foreach\y in{.1,.3,.5}{\draw[guide](0,\y)--({4*\y},{5*\y});}\foreach\y in{1.3,2.1,2.9,3.7,4.5}{\draw[guide](0,\y)--({(4-.8*\y)/1.8},{(4+\y)/1.8});}}
\end{exercise}
\begin{solution}
去程接收者远离，频率比为$\sqrt{(1-\beta)/(1+\beta)}=1/3$；返程靠近，频率比为3。旅行者两段固有时间各7.5年，所以去程接收相当于地球2.5年内发出的信号，返程接收其余22.5年。
\Diagram{fig:m2signalsa}{转向瞬间接收到的光来自地球2.5年时刻}{\draw[axis](-.2,0)--(3.5,0)node[right]{$x$};\draw[axis](0,-.2)--(0,5.5)node[above]{$ct$};\draw[curve](0,0)--(2,2.5)--(0,5);\draw[very thick,->](0,.5)--(2,2.5);\node[left]at(0,.5){2.5年};\node[right]at(2,2.5){$B$};}
图\ref{fig:m2signalsa}中由光速连线也可算发射时刻$t_e=t_B-l/c=12.5-10=2.5$年。与旅行者去程、返程赋予地球的同时刻4.5年、20.5年都不同。接收信号连续，只有接收频率在理想瞬时换向模型中突变。
\end{solution}

\begin{exercise}\label{ex:m2fuelreturn}
光子火箭需由地球静止起飞，加速至$\beta c$，到远处减速至相对地球静止，随后原路返回并最终停下。忽略巡航耗能，求包括全部返航燃料在内的初末质量比。
\Diagram{fig:m2returnq}{四段速度改变与两段巡航}{\draw[axis](-.3,0)--(6,0)node[right]{$x$};\fill(0,0)circle(1.5pt)node[below]{地球};\fill(5,0)circle(1.5pt)node[below]{远端};\draw[->](.3,.8)--(4.7,.8)node[midway,above]{$+\beta c$};\draw[->](4.7,-.8)--(.3,-.8)node[midway,below]{$-\beta c$};}
\end{exercise}
\begin{solution}
令$\chi=\operatorname{artanh}\beta>0$。图\ref{fig:m2returnq}中每次由静止到巡航或由巡航到静止，都需消耗绝对快度变化$\chi$。喷流方向在减速时反转，火箭质量仍减少，故应对快度变化取绝对值。
\Diagram{fig:m2returna}{快度变化总长度为$4\chi$}{\draw[axis](-.3,0)--(6,0)node[right]{过程};\draw[axis](0,-1.5)--(0,1.8)node[above]{$\chi_{\rm ship}$};\draw[very thick](0,0)--(1,1)--(2,1)--(3,0)--(4,-1)--(5,-1)--(6,0);\node[left]at(0,1){$\chi$};\node[left]at(0,-1){$-\chi$};}
图\ref{fig:m2returna}中四次点火的质量比逐次相乘，故$M_0/M_f=\ee^{4\chi}=[(1+\beta)/(1-\beta)]^2$。若$\beta=.8$，质量比为81。不能先算单程消耗，再以最终干质量为基准简单加倍，因为起飞时必须推动全部后续燃料。
\end{solution}

\begin{exercise}\label{ex:m2constalpha}
飞船从静止出发，沿$x$保持固有加速度$\alpha$，直到到达地球距离$L$处。求地球时间、船上时间和末速度。
\Diagram{fig:m2alphaq}{固定终点距离的持续加速}{\draw[axis](-.3,0)--(5,0)node[right]{$x$};\fill(0,0)circle(1.5pt);\draw[curve](4,-.5)--(4,1.2);\draw[<->](0,-.5)--(4,-.5)node[midway,below]{$L$};\draw[axis](.5,.5)--(2,.5)node[above]{$\alpha$};}
\end{exercise}
\begin{solution}
图\ref{fig:m2alphaq}取飞船由原点静止出发。固有加速度定义给$\dd\chi/\dd\tau=\alpha/c$，所以$\chi=\alpha\tau/c$。再由四速度的时间和空间分量，
\[
 t=\int_0^\tau\cosh(\alpha s/c)\dd s
 =\left[\frac c\alpha\sinh(\alpha s/c)\right]_0^\tau
 =\frac c\alpha\sinh\chi,
\]
\[
 x=\int_0^\tau c\sinh(\alpha s/c)\dd s
 =\left[\frac{c^2}\alpha\cosh(\alpha s/c)\right]_0^\tau
 =\frac{c^2}\alpha(\cosh\chi-1).
\]
两个积分常数已由$t=x=0$确定。
\Diagram{fig:m2alphaa}{终点落在恒固有加速度双曲线上}{\draw[axis](-.2,0)--(4.6,0)node[right]{$\alpha x/c^2$};\draw[axis](0,-.2)--(0,4.6)node[above]{$\alpha t/c$};\draw[curve,domain=0:2,samples=70]plot({cosh(\x)-1},{sinh(\x)});\fill(2,2.828)circle(1.5pt)node[right]{$x=L$};\draw[guide](2,0)--(2,2.828);}
在图\ref{fig:m2alphaa}中令$q=\alpha L/c^2$。到达时$\cosh\chi_f=1+q$，正向运动选择$\sinh\chi_f=+\sqrt{(1+q)^2-1}=\sqrt{q(2+q)}$，因此
\[
 t_f=\frac c\alpha\sqrt{q(2+q)},\quad
 \tau_f=\frac c\alpha\operatorname{arcosh}(1+q),\quad
 \beta_f=\frac{\sqrt{q(2+q)}}{1+q}.
\]
平方后$1-\beta_f^2=(1+q)^{-2}>0$，也给$\gamma_f=1+q$。小$q$时$\operatorname{arcosh}(1+q)\simeq\sqrt{2q}$，两种时间同趋$\sqrt{2L/\alpha}$、末速趋$\sqrt{2\alpha L}$。大$q$时$t_f\simeq L/c+c/\alpha$，而$\tau_f\simeq(c/\alpha)\ln(2q)$，远距离的船上时间仅呈对数增长；这依赖全程持续提供恒固有加速度。
\end{solution}

\begin{exercise}\label{ex:m2powerengine}
理想光子火箭在自身瞬时静止系的辐射功率恒为$P_*$。将它与恒固有加速度光子火箭比较：分别求质量随固有时间的变化及固有功率随时间的条件。
\Diagram{fig:m2engineq}{定向光束的固有功率与推力}{\draw[curve](2,0)rectangle(4.3,1);\node at(3.15,.5){火箭$m$};\draw[axis](1.8,.5)--(.2,.5)node[above]{$P_*/c$};\draw[axis](4.5,.5)--(5.7,.5)node[above]{$\alpha$};}
\end{exercise}
\begin{solution}
图\ref{fig:m2engineq}中，瞬时船系向后发光的能量率为$P_*$，动量率大小为$P_*/c$。一阶能量与动量守恒分别给
\[
 \frac{\dd m}{\dd\tau}=-\frac{P_*}{c^2},\quad
 mc\frac{\dd\chi}{\dd\tau}=\frac{P_*}{c},\quad
 \alpha=c\frac{\dd\chi}{\dd\tau}=\frac{P_*}{mc}.
\]
若$P_*$恒定，以$m(0)=M_0$积分：$\int_{M_0}^{m}\dd\mu=-(P_*/c^2)\int_0^\tau\dd s$，得到$m=M_0-P_*\tau/c^2$。再相除两个微分方程，有$\dd\chi=-\dd m/m$，从初始静止积分得$\chi=-[\ln\mu]_{M_0}^{m}=\ln(M_0/m)$。推力恒定而惯性质量减小，故固有加速度不断增大。
\Diagram{fig:m2enginea}{恒加速度需要功率与剩余质量同步下降}{\draw[axis](0,0)--(4.3,0)node[right]{$\alpha\tau/c$};\draw[axis](0,0)--(0,2.8)node[above]{$m/M_0,\ P_*/P_{*0}$};\draw[curve,domain=0:3.5,samples=70]plot(\x,{2.3*exp(-\x)});\node[left]at(0,2.3){$1$};}
若要求图\ref{fig:m2enginea}中的$\alpha$恒定，必须令功率随质量改变：$P_*=mc\alpha$。质量方程于是变为
\[
 \frac{\dd m}{m}=-\frac\alpha c\dd\tau,
 \quad \int_{M_0}^{m}\frac{\dd\mu}{\mu}
 =-\frac\alpha c\int_0^\tau\dd s,
 \quad \ln\frac m{M_0}=-\frac{\alpha\tau}c.
\]
取指数并代回功率条件，
\[
 m=M_0\ee^{-\alpha\tau/c},\qquad
 P_*=M_0c\alpha\ee^{-\alpha\tau/c}.
\]
此时$\chi=\alpha\tau/c$，与恒固有加速度世界线一致。若干质量为$M_d$，两种程序都在$m=M_d$时停机：恒功率的可燃固有时间为$(M_0-M_d)c^2/P_*$，恒加速度则为$(c/\alpha)\ln(M_0/M_d)$。在相同初末质量比下，两种程序给出相同末快度；燃烧历时和沿途位置由功率程序决定。
\end{solution}

\begin{exercise}\label{ex:m2oppositejets}
一对同时由同一点射出的对称光斑在源系以同速$\beta c$沿相反方向运动。视线与靠近光斑速度夹角为$\theta$。测得两侧视横向速度分别为$u_+,u_-$。求$\beta\cos\theta$及$\beta\sin\theta$。
\Diagram{fig:m2jetsq}{靠近和远离光斑的投影}{\draw[axis](-.3,0)--(4.5,0)node[right]{视线};\draw[very thick,->](0,0)--(2.7,1.7)node[right]{$+$};\draw[very thick,->](0,0)--(-1.9,-1.2)node[left]{$-$};\node[below right]at(0,0){源};}
\end{exercise}
\begin{solution}
图\ref{fig:m2jetsq}的两侧横向位移大小相同，但接收时间分别缩短和拉长，故
\[
 \frac{u_+}c=\frac{\beta\sin\theta}{1-\beta\cos\theta},\quad
 \frac{u_-}c=\frac{\beta\sin\theta}{1+\beta\cos\theta}.
\]
当$u_++u_->0$时，两式相除后解得$A=\beta\cos\theta=(u_+-u_-)/(u_++u_-)$。再乘回任一分母，得$B=\beta\sin\theta=2u_+u_-/[c(u_++u_-)]$。若两侧视横向速度都为零，可能是$\beta=0$，也可能是运动轴沿视线；此时横向数据不足以唯一反演，不能使用上述零分母公式。
\Diagram{fig:m2jetsa}{两个可测投影组合给出速度和视角}{\draw[axis](0,0)--(3.8,0)node[right]{$A$};\draw[axis](0,0)--(0,2.8)node[above]{$B$};\draw[very thick,->](0,0)--(2.7,1.8)node[right]{$\beta$};\draw[guide](2.7,0)--(2.7,1.8)--(0,1.8);\node at(.9,.3){$\theta$};}
图\ref{fig:m2jetsa}给$\beta=\sqrt{A^2+B^2}$、$\theta=\operatorname{atan2}(B,A)$。若由观测反演得到$\beta\geq1$，则对称、同速、远场或距离标定中至少一项不满足模型假设，不能把不相容的数据解释成真实超光速。
\end{solution}

\begin{exercise}\label{ex:m2imagebar}
长度$L_0$的细杆沿$x$运动，杆与速度平行。观察者在远处，视线单位矢量$\vct n$与$x$夹角$\theta$。求同一照片中杆端横向投影的长度，并与等时长度比较。
\Diagram{fig:m2barq}{照片由两端不同的发光事件组成}{\draw[curve](0,0)--(3.5,0);\fill(0,0)circle(1.5pt)node[below]{$A$};\fill(3.5,0)circle(1.5pt)node[below]{$B$};\draw[axis](1.7,.5)--(3,.5)node[above]{$V$};\draw[guide](0,0)--(1.5,2.2);\draw[guide](3.5,0)--(5,2.2);\node[right]at(5,2.2){$\vct n$};}
\end{exercise}
\begin{solution}
图\ref{fig:m2barq}中，实验室同时杆长为$L=L_0/\gamma$。取尾端、头端世界线$x_A=Vt_A$、$x_B=Vt_B+L$。远处观察者的方向$\vct n$由物体指向相机，源点到相机的距离差近似为$-\vct n\cdot\Delta\vct r_e=-\cos\theta\Delta x_e$。同一照片要求
\[
 0=\Delta t_o=t_B-t_A-\frac{\cos\theta}{c}\Delta x_e.
\]
再由两端世界线相减，$\Delta x_e=V(t_B-t_A)+L$，消去发射时间差得
\[
 \Delta x_e=\frac{L}{1-\beta\cos\theta}.
\]
\Diagram{fig:m2bara}{发光事件空间差向成像面的投影}{\draw[curve](0,0)--(4,0);\draw[guide](0,0)--(2,2.5);\draw[guide](4,0)--(5.56,1.95);\draw[very thick](2,2.5)--(4.44,.55);\draw[guide](4,0)--(4.44,.55);\node[above right]at(3.4,1.6){$L_{\rm photo}$};\node[below]at(2,0){$\Delta x_e$};}
在图\ref{fig:m2bara}的像平面内，端点分离矢量须扣去沿视线分量，故
\begin{align*}
 L_{\rm photo}^2
 &=|\Delta\vct r_e|^2-(\vct n\cdot\Delta\vct r_e)^2\\
 &=(\Delta x_e)^2\sin^2\theta,
\end{align*}
因此$L_{\rm photo}=L_0\sin\theta/[\gamma(1-\beta\cos\theta)]$。在$\theta=\pi/2$时恢复$L_0/\gamma$。其角度函数与例\ref{ex:m2super}相同，沿用该例已经展开的商式求导，最大值在$\cos\theta=\beta$取得；代入$\sin\theta=1/\gamma$，得$L_{\rm photo,max}=L_0$。式中的$L$、$\Delta x_e$、$L_{\rm photo}$分别表示同时长度、延迟发射位置差和像面投影长，三者须按各自定义使用。
\end{solution}

